
\subsection{A Verifiable Random Function without Pairings}

We now describe a verifiable random function that is pairing-free.
We build upon the VRF defined by Dodis-Yampolskiy VRF~\cite{DodisY05},
but where the proving mechanism is performed using discrete logarithm proofs of equality (DLEQ),
which are pairing-free.

We first begin by defining verifiable random functions more generally.
A verifiable random function (VRF) is a keyed PRF,
such that the PRF can be evaluated using only knowledge of the secret key,
but is publicly verifiable given the public key.

\begin{definition}
A \defn{verifiable random function} (VRF) $\VRF$ is a tuple of algorithms
$\VRF = (\setup, \keygen, \prove, \verify)$, such that:

\begin{itemize}
    \item $\VRF.\setup(1^\lambda) \rightarrow \pp$:
Accepts as input a security
parameter $\lambda$ and outputs public parameters $\pp$,
which are then implicitly given as input to all other algorithms.
    \item $\VRF.\keygen() \rightarrow (\sk, \pk)$:
A probabilistic algorithm that outputs a VRF secret key $\sk$ and public key $\PK$.
      \item $\VRF.\prove(\sk, \vrfin) \rightarrow (\vrfout, \pi)$:
Accepts as input a secret $\sk$ and value $\vrfin$ to evaluate the VRF on.
Outputs the evaluation of $x$ as $\vrfout$ and a proof $\pi$ that the evaluation was performed correctly.
      \item $\VRF.\verify(\pk, \vrfin, \vrfout, \pi) \rightarrow \zo$:
Accepts as input the input to evaluate $\vrfin$,
a public key $\pk$,
the VRF output $\vrfout$,
and a proof $\pi$.
Outputs $1$ if $\vrfout$ is a valid output with respect to $\pk$ and $\vrfin$,
otherwise, outputs $0$.
\end{itemize}
\end{definition}

\paragraph{Correctness.}
A VRF is \defn{correct} if for all $\lambda$ and $\pp \gets \VRF.\setup(1^\lambda)$
and $(\PK, \PK) \rgets \VRF.\keygen()$,
Equation~\ref{eq:vrf-correct} holds:

\begin{align} \label{eq:vrf-correct}
  \begin{split}
    \VRF.\verify(\PK, \vrfin, \pi) \rightarrow 1 \text{ when} \\
    \VRF.\prove(\SK, \vrfin) \rightarrow (\vrfout, \pi)
  \end{split}
\end{align}

\paragraph{Security.}
A VRF is \emph{secure} if it is unique and pseudorandom.
We define these notions next.

\[
  \advant{unq}{\adv,\VRF}(\lambda)
     = \Pr\left[
            \begin{array}{l}
                \pp \gets \VRF.\Setup(1^\lambda) \\
               (\PK, \vrfin,\vrfout, \pi, \vrfout', \pi') \rgets \adv(\pp)
            \end{array}
            \ \ :\ \
            \begin{array}{l}
                \VRF.\verify(\PK, \vrfin, \vrfout, \pi) = 1 \\
                \VRF.\verify(\PK, \vrfin, \vrfout', \pi') = 1 \\
                \vrfout \neq \vrfout'
            \end{array}
        \right]
\]

\begin{definition}
  A VRF is \defn{unique} if for all PPT adversaries $\adv$,
  $\advant{unq}{\adv,\VRF}$ is a negligible function of $\lambda$.
\end{definition}


\[
  \advant{psdr}{\adv,\VRF}(\lambda)
     = \Pr\left[
            \begin{array}{l}
              b = b', \text{ and} \\
              \oracle^{\prove} \text{ was not queried on } \vrfin
            \end{array}
            \ \ :\ \
            \begin{array}{l}
                b \rgets \zo \\
                \pp \gets \VRF.\Setup(1^\lambda) \\
                (\sk, \pk) \rgets \VRF.\keygen() \\
                \vrfin \rgets \adv(\pp, \pk) \\
                (\vrfout_0, \pi) \rgets \VRF.\prove(\SK, \vrfin) \\
                \vrfout_1 \rgets \vrfcodomain \\
                b' \rgets \adv^{\prove}(\vrfout_b, \pstate_A)
            \end{array}
        \right]
\]

\begin{definition}
  A VRF is \defn{pseudorandom} if for all PPT adversaries $\adv$,
  $\advant{psdr}{\adv,\VRF}$ is a negligible function of $\lambda$.
\end{definition}


\subsubsection{A Verifiable Random Function without Pairings}
\label{sec:dyprf}

We now give a variant of the Dodis-Yampolskiy VRF~\cite{DodisY05},
which itself builds on the Boneh-Boyen-Shacham unpredictable function~\cite{BonehBS04}.
However, we define this VRF without relying upon pairings,
by instead employing discrete log equality proofs,
as further described in Section~\ref{sec:dleq}.
We refer to this variant of the Dodis-Yampolskiy VRF as $\DYPRF$.

$\DYPRF$ is defined respect to the hash function $\Hash$,
such that:

  \begin{itemize}
    \item $\DYPRF[\Hash].\setup(1^\lambda) \rightarrow \pp$:
      Perform $(\GG, q, g) \gets \GroupGen$.
      Then, sample $\omega \rgets \Zq$ and derive $h \gets g^\omega$.
      Output $(\GG, q, g, h) $.
    \item $\DYPRF[\Hash].\keygen() \rightarrow (\sk, \pk)$:
        Sample $\sk \rgets \Zq$; derive $\pk \gets g^{\sk}$.
        Output $(\sk, \pk)$.
      \item $\DYPRF[\Hash].\prove(\sk, x) \rightarrow (Y, \pi)$:
        On input the secret key $\sk$ and value $x \in \Zq$,
        derive $\alpha \gets \sk + x$.
        derive $ Y \gets g^{1/\alpha}$.
        Then, obtain $\pi \gets \DLEQ[\Hash].\prove(g, g^\alpha, Y, g, \alpha)$.
        Output $(Y, \pi)$.
      \item $\DYPRF[\Hash].\verify(x, \pk, Y, \pi) \rightarrow \zo$:
        Derive $B \gets \PK \cdot g^x$.
        Then, perform the following check: $\DLEQ[\Hash].\verify(g, B, Y, g, \pi) \iseq 1$.
        Output $1$ if the check holds, otherwise, output $0$.
  \end{itemize}

\paragraph{Correctness.}
$\DYPRF$ is correct, because the proof $\pi$ will be of the form $\pi = (R_1, R_2, z)$,
where $R_1 = g^r, R_2 = g^{\frac{r}{\SK + \vrfin}}$,
and $z = r + c(\SK + \vrfin)$.
And so the following relations in Equations~\ref{eq:correct-one} and~\ref{eq:correct-two} will hold:

\begin{align}
  \label{eq:correct-one}
  \begin{split}
    g^z &= R_1 \cdot B^c \\
    g^{r + c(\SK + \vrfin)}  &= g^r \cdot (g^{\SK + \vrfin})^c
  \end{split}
\end{align}

\begin{align}
  \label{eq:correct-two}
  \begin{split}
    Y^z &= R_2 \cdot g^c \\
    g^{\frac{r}{\SK + \vrfin} + \frac{c(\SK + \vrfin)}{\SK + \vrfin}}  &= g^\frac{r}{\SK + \vrfin} \cdot g^c \\
    g^{\frac{r}{\SK + \vrfin} + c}  &= g^\frac{r}{\SK + \vrfin} \cdot g^c
  \end{split}
\end{align}

\paragraph{Security.}
Recall that for a VRF to be secure,
it must be unique, verifiable, and pseudorandom~\cite{DodisY05}.
The pseudorandomness of $\DYPRF$ in the non-pairing setting is given by Dodis and Yampolskiy~\cite{DodisY05},
and assumes the Diffie-Hellman inversion (DHI) assumption.
Because DLEQ proofs are complete and sound,
then $\DYPRF$ is also unique and verifiable.

\subsection{Chaum-Pedersen Discrete-Log Equality Proofs}
\label{sec:dleq}

We now describe how to generate proofs of discrete log equality between two
group elements,
via Chaum-Pedersen proofs~\cite{BonehS2020}.

\begin{definition}
A Chaum-Pedersen discrete-log equality proof $\DLEQ$ is the tuple of algorithms
  $(\DLEQ = (\setup, \prove, \verify))$,
where:

\begin{itemize}
  \item $\DLEQ[\Hash].\setup(1^\lambda) \rightarrow \pp$:
Derive $(\GG, q, g) \gets \GroupGen(1^\lambda)$.
Sample $h \rgets \GG$.
Output $\pp = (\GG, q, g, h)$.
  \item $\DLEQ[\Hash].\prove(A, B, C, D, w) \rightarrow \pi$:
Accepts as input four group elements $(A, B, C, D) \in \GG^4$ and a value $w \in \Zq$,
with the following relation:

    \begin{align}
      \begin{split}
        B &= A^w, \text{ and} \\
        D &= C^w
      \end{split}
    \end{align}

Performs the following steps:
\begin{enumerate}
\item Sample $r \rgets \Zq$ and derive $R_1 \gets A^r$ and $R_2 \gets C^r$.
\item Derive $c \gets \Hash(A, B, C, D, R_1, R_2)$.
\item Then, derive $z \gets r + c \cdot w$.
\item Output $\pi = (R_1, R_2, z)$.
\end{enumerate}
\item $\DLEQ[\Hash].\verify(A, B, C, D, \pi) \rightarrow \zo$
Accepts as input four group elements $(A, B, C, D) \in \GG^4$,
and a proof $\pi$.
Performs the following steps:
\begin{enumerate}
\item Parse $(R_1, R_2, z) \gets \pi$.
\item Then, derive $c' \gets \Hash(A, B, C, D, R_1, R_2)$.
\item Finally, perform the checks in Equation~\ref{eq:dleq-verify}:
  \begin{align} \label{eq:dleq-verify}
\begin{split}
  & A^z \iseq R_1 \cdot B^c, \text{ and} \\
  & C^z \iseq R_2 \cdot D^c
\end{split}
\end{align}
\item Output $1$ if both checks hold, otherwise, output $0$.
\end{enumerate}
\end{itemize}
\end{definition}

\paragraph{Security.}
DLEQ proofs are special-sound and honest-verifier zero-knowledge~\cite{BonehS2020}.


